% Generated draft: requires mathematical and source review.
\documentclass[11pt]{article}
\usepackage[margin=25mm]{geometry}
\usepackage{fontspec}
\setmainfont{Latin Modern Roman}
\usepackage{amsmath,amssymb,mathtools,unicode-math}
\setmathfont{Latin Modern Math}
\usepackage{enumitem,needspace,xurl,hyperref}
\setlength{\parindent}{0pt}
\setlength{\parskip}{5pt}
\setlength{\emergencystretch}{3em}
\newcommand{\N}{\mathbb{N}}
\newcommand{\Z}{\mathbb{Z}}
\newcommand{\Q}{\mathbb{Q}}
\newcommand{\R}{\mathbb{R}}
\newcommand{\C}{\mathbb{C}}
\newcommand{\F}{\mathbb{F}}
\newcommand{\A}{\mathbb{A}}
\newcommand{\PP}{\mathbb{P}}
\newcommand{\E}{\mathbb{E}}
\newcommand{\eps}{\varepsilon}
\newcommand{\dd}{\,\mathrm{d}}
\begin{document}
\section*{UnsolvedMath ID: 20000249}
\textbf{Source problem code:} catalogue entry
\subsection{Definitions and mathematical statement}
Anastasia Stavrova's AIM question asks whether \(\pi_0^{\A^1}(G)\) is always abelian for the intended class of linear algebraic groups, and whether the known case-by-case abelianity results admit a uniform proof. The OCR wording ``linear reduced/algebraic group'' does not fix the precise class of groups or base-field hypotheses. For a Nisnevich sheaf of groups \(\mathcal H\) on smooth \(k\)-schemes, \(\pi_0^{\A^1}(\mathcal H)\) means the sheaf of genuine \(\A^1\)-homotopy components; abelianity means vanishing of its commutator morphism, equivalently abelian groups of sections on every smooth test scheme.

The supplied report proposes the arbitrary-base-field criterion
\[
\pi_0^{\A^1}(\mathcal H)\text{ is abelian}
\quad\Longleftrightarrow\quad
\pi_0^{\A^1}(\mathcal H)(K)\text{ is abelian for every finitely generated separable field extension }K/k.
\]
Here ``separable'' includes the usual separably generated function fields, not merely finite algebraic extensions; field values use the standard extension from smooth models. The report attributes the reduction to known generic-point injectivity applied to commutators, but does not give its precise theorem or reference.

For semisimple simply connected algebraic groups it invokes a fieldwise comparison
\[
\pi_0^{\A^1}(G)(K)\cong G(K)/R,
\]
and calls the abelianity question a uniform reduction to the long-open abelianity problem for \(R\)-equivalence. Here \(R\)-equivalence is generated by rational curves in \(G\) defined at the two specified endpoints; those curves need not extend to regular maps on the entire affine line. The supplied summary claims the comparison unconditionally in that scope but omits precise hypotheses and bibliographic identification. This reduction is not a solution to all abelianity cases.

The original remark of Madhav Nori says the assertion is not always true, proposes \(SU(2,\R)\), and suggests the same reasoning for any anisotropic group over \(\R\). The attached report instead interprets that example as the norm-one quaternion group and proposes, for every real closed field \(F\),
\[
G=SL_1((-1,-1)_F),\qquad
\text{every regular map }\A^1_F\to G\text{ is constant},
\]
\[
\mathcal S(G)(F)=G(F)\text{ is nonabelian},\qquad
G(F)/R=1,\qquad \pi_0^{\A^1}(G)(F)=1.
\]
Here \(\mathcal S\) denotes naive affine-line chain components, and \((-1,-1)_F\) is the quaternion algebra. The report invokes Cayley rational curves for the \(R\)-equivalence assertion. Thus it presents the workshop annotation as confusing naive chain components with genuine components. This is a reported correction requiring verification, not an unqualified repetition of the annotation. Nor can the literal phrase ``any anisotropic group'' imply nonabelian components for anisotropic commutative groups; the remark's intended scope needs qualification. Triviality at real closed fields alone is not asserted to identify the entire component sheaf over all field extensions.
\subsection{Short English statement}
Ask when genuine \(\A^1\)-component sheaves of algebraic groups are abelian. The report proposes a fieldwise test and an \(R\)-equivalence reduction, and distinguishes a nonabelian naive quaternion example from its claimed trivial genuine components over real closed fields.
\subsection{Sources}
\begin{itemize}
\item AIM workshop ``Projective modules and \(A^1\)-homotopy theory,'' May 5--9, 2014, Question 12, attributed to Anastasia Stavrova, with a remark by Madhav Nori: \url{https://aimath.org/pastworkshops/projectiveA1problems.pdf}. The verification field claims official-PDF confirmation with modernized typography; no PDF was accessed here.
\item Local canonical registry \texttt{aim-algebraic-geometry-notes.json}, \texttt{notes[248]}, tag \texttt{question}; permanent UnsolvedMath ID \(20000249\), source code \texttt{catalogue entry}. Original asks about abelianity and a uniform argument; the nested remark proposes \(SU(2,\R)\) and anisotropic real groups as counterexamples. The original literature field is empty.
\item Supplied research attempt 1 proposing the generic-point criterion and real-closed quaternion correction. No full authoritative references for generic-point injectivity, the fieldwise comparison, or the candidate correction are supplied. The report is machine generated, not peer reviewed, and requires independent verification before citation.
\end{itemize}
\end{document}
